%----------
%	CLASIFICACIÓN ZERO-SHOT CON LIBRO DE CÓDIGOS
%----------

\label{sec:zeroshot_iris}

La clasificación se aborda en régimen \textit{zero-shot}, presentado en la
Sección~\ref{sec:ai_solution}. Al no disponerse de un corpus de entrenamiento
etiquetado a gran escala, la supervisión por ejemplos se sustituye por la
\textbf{supervisión mediante libro de códigos}~\cite{ziems2024css, hollauer_2024},
de modo que el modelo generaliza a partir del conocimiento que ya posee y de la
metodología que recibe en la instrucción.

La metodología que se traslada al \textit{prompt} procede de la especificación
estructurada de cada variable descrita en la Sección~\ref{sec:iris_variables}:
criterio de decisión, pasos, ejemplos y fronteras. Todo ese conocimiento experto
se centraliza en un único fichero, \texttt{variables.json},
que actúa como fuente de verdad del sistema. A partir de él se generan de forma
automática las dos arquitecturas que se comparan en el
Capítulo~\ref{implementation}, de modo que ambas parten de idéntico contenido y
difieren solo en cómo lo reciben.

\subsubsection{Salida estructurada y trazabilidad}

Por cada pieza y variable, el modelo devuelve un JSON con tres campos:

\begin{itemize}
    \item \texttt{codigo}: el valor en la escala cruda de la variable
    (1 = ausencia, 2 = presencia, y V25 admite además el 3 para el salto semántico).
    El modelo emite este código sin agrupar. Solo al calcular las métricas los
    valores 2 y 3 se unen como clase positiva (Sección~\ref{sec:iris_variables}).
    \item \texttt{explicacion}: la justificación, que debe citar los pasos de la
    metodología aplicados.
    \item \texttt{evidencias}: fragmentos literales del texto que sostienen la
    decisión (lista vacía cuando \texttt{codigo}\,=\,1, es decir, en los negativos).
\end{itemize}

\begin{figure}[H]
    \centering
    \includegraphics[width=\textwidth]{Plantilla_TFG_ingles_2019/imagenes/salida_trazabilidad_blanco.pdf}
    \caption{Salida estructurada y trazable de una predicción.}
    \label{fig:iris-salida}
\end{figure}

Este formato materializa el marco
\textit{human-in-the-loop}. La explicación y las evidencias hacen la predicción
auditable: la analista lee la justificación y decide si acepta, rechaza o matiza
el veredicto. Detectar un sesgo de género no se reduce a localizar palabras
clave, sino que exige interpretar el contexto social del texto, un criterio que
sigue siendo competencia humana. Por eso el sistema asiste a la analista en lugar
de sustituirla.

\subsubsection{Ventajas y límites del enfoque}

El alcance de esta decisión se aprecia mejor si se examinan sus efectos sobre dos
planos distintos, el de la adquisición de conocimiento y el del modelado
estadístico de la tarea.

En el plano de la \textbf{adquisición de conocimiento}, el sistema no incorpora
nada de lo que procesa. Los pesos del modelo permanecen intactos, de modo que la
metodología no llega a formar parte de él, sino que se le suministra en cada
llamada y se pierde al terminarla. Esto tiene una cara favorable y otra adversa.
La favorable es que prescinde de conjuntos de entrenamiento etiquetados y del
preprocesamiento intensivo de los enfoques clásicos, y sobre todo que mejorar el
sistema consiste en reescribir la metodología y no en reentrenar un modelo, tarea
que corresponde entonces al equipo experto y no a quien programa. La adversa es
que el sistema no acumula experiencia, ya que un error cometido en una pieza no
informa el análisis de la siguiente, y todo el conocimiento del dominio con que
cuenta procede de su preentrenamiento, ajeno a este corpus y a esta metodología.

En el plano del \textbf{modelado estadístico}, el sistema no estima probabilidad
alguna a partir de los datos anotados. Un clasificador supervisado aprendería con
qué frecuencia aparece cada clase y ajustaría en consecuencia su umbral de
decisión, de manera que en una variable como V33, con un 6,2\,\% de positivos,
tendería por sí solo a marcar poco. Aquí no sucede nada de eso. Al no observar
nunca la distribución de clases del corpus (Sección~\ref{sec:iris_variables}), la
propensión del modelo a marcar no responde a la prevalencia real, sino a cómo
interpreta el criterio que se le ha descrito y a las regularidades de su
preentrenamiento. Cuando marca de más o de menos no está, por tanto, mal
calibrado respecto a una distribución que conoce, sino aplicando un umbral
interpretativo propio, observación que resulta determinante para leer los
resultados del Capítulo~\ref{results_and_discussion} y, en particular, las
diferencias de recall entre modelos.

De ambos planos se sigue la limitación principal del enfoque. El rendimiento queda
determinado por la forma en que la metodología se traslada al \textit{prompt} y
por la capacidad del modelo empleado, sin que existan ejemplos etiquetados con los
que corregir sus desviaciones. Determinar el alcance de esa limitación, y
establecer si su origen reside en el modelo, en la forma de presentarle la
metodología o en la definición misma de la tarea, constituye el objeto del
Capítulo~\ref{results_and_discussion}.
